% =====================================================================
%  Chapter 17 — Lightcurve Inversion\idx{Lightcurve inversion} and the Convex Hull
% =====================================================================

\chapter{Lightcurve Inversion and the Convex Hull}
\label{ch:lightcurve}

The lightcurve inversion\idx{Lightcurve inversion} is the workhorse of asteroid characterisation. It
has produced more shape models than every other method combined, it uses
telescopes that already exist, and it returns a coherent three-dimensional
body with a spin axis and a rotation period. This chapter asks what it
returns and what the answer costs.

\section{What the inversion does}

The procedure is standard and worth stating in full, because the place the
assumption enters is usually a single sentence that is easy to miss.

A set of photometric observations of a rotating body is recorded. A model is
constructed: a shape, an orientation and a rotation rate, such that the
model's predicted brightness variation as it rotates matches the observed
variation. The model parameters are then searched for the best fit.

The fitted parameters are the rotation period, the direction of the spin
pole, and the shape. The shape that comes out of the standard inversion is
the \emph{convex hull} of the body consistent with the observations ---
because the constraint set is not convex and the hull is the standard way to
report a single answer from it.

The assumption is not a modelling convenience. It is the statement that the
body is convex, and it is made silently, because a convex body is the
simplest consistent answer and because the database that holds these models
records it.

\section{What the record contains}

The scale of the effort is remarkable and it is the reason the method
dominates the field. The asteroid models database\idx{Asteroid models database} reached roughly 25,000
models by October 2026, a figure that tripled the catalogue in five years.
The sample underlying the most recent analysis covered 881 near-Earth
asteroids, of which 47 received a model.

The throughput comparison that comes with it is the most useful summary of
the field's measurement history there is. Taken across the epochs of the
field's major reviews, the number of resolved shape models produced by each
method has moved as follows:

\begin{table}[!htbp]
\centering
\small
\caption{Resolved shape models by method, across the epochs of the field's
major reviews. The figures are the accumulated totals.}
\label{tab:throughput}
\begin{tabularx}{\textwidth}{@{}P{1.15in}P{0.78in}P{0.78in}Y@{}}
\toprule
\rowcolor{panel}
\textbf{Method} & \textbf{Earliest} & \textbf{Latest} & \textbf{Comment} \\
\midrule
Interferometry & 5 & \(10^{2}\) & the only resolved method with a
projected hundred \\
\rowcolor{panel}
Photometry & 30 & \(10^{4}\) & the workhorse; lightcurve inversions \\
Fly-by & 6 & fewer than 15 & one per spacecraft encounter \\
\bottomrule
\end{tabularx}
\end{table}

\tabref{tab:throughput} contains the argument of this book in one line.
Interferometry, the only method that produces a resolved shape without
requiring proximity, is the lowest-yield resolved method in the field, and
the field's own roadmap projects a hundred of them. The ten thousand
photometric inversions are low-order convex bodies. The fewer than fifteen
fly-by images are the only true three-dimensional shapes in existence.

\section{The cost of the convexity, and the evidence for it}

The question the convex-hull assumption raises is quantitative: how wrong is
a convex hull, for a real small body?

The best available evidence is the one measurement that bypassed the
question, and it is the 2017 radio occultation of (372) Palma discussed in
Chapter 8. The silhouette it measured was slightly elongated and had a
concave edge, and the roundness\idx{Roundness} parameter --- the ratio of the inscribed to
the circumscribed circle --- had a mean of 0.866 with a standard deviation of
0.093 over the fitted shape ensemble, giving a departure from a perfect
circle of ten to twenty per cent.

Twenty per cent is not a small number, and it is worth being precise about
what it does and does not imply. The roundness\idx{Roundness} parameter is a two-dimensional
measure of the silhouette. It is not the volume error. A body that is
slightly elongated in projection may be considerably less voluminous than
its convex hull if it is also concave, and considerably more voluminous than
its projection if it is oblate in the third dimension. The honest statement
is a bound and a direction: \textbf{real small bodies depart from the
convex ideal by tens of per cent in their silhouettes, so convex-hull volumes
carry systematic errors of that order, and the errors are one-sided --- the
hull is always at least as large as the body.}

\begin{seamnote}{Why the one-sidedness is good news}
Because the convex hull always overstates the volume, and therefore
overstates the mass at a fixed density, a lightcurve inversion\idx{Lightcurve inversion} tends to
overstate the energy. That is a conservative error in the direction of
concern. It is also an error that nobody has bounded, because the field does
not have the resolved silhouettes of small bodies against which to bound it.
A single measurement that would bound it exists, and it is the measurement
the rest of this book is about.
\end{seamnote}

\section{The second limitation, which is not about shape}

The convex hull is the assumption that gets discussed. The limitation that
gets less attention is that a lightcurve inversion returns a shape, and a
shape is not a size.

The diameter that goes into the energy chain is not the diameter of the
shape model. It is the diameter consistent with the model's projected area
and the object's absolute magnitude --- which is a photometric inversion with
its own assumptions, and which carries the albedo into the answer. The
lightcurve inversion therefore delivers a shape, and then a separate step
turns the shape into a diameter, and that step is where the albedo
degeneracy re-enters.

A resolved image does not have this problem in the same way, because the
projected area is measured directly and the albedo is an output rather than
an input. That difference is the entire content of the argument in
\tabref{tab:routes}.

\section{What the method is good for}

It is worth ending this chapter by saying plainly what the lightcurve
inversion is good for, because the chapter is not a criticism of it.

It is good for the spin state. The period and the pole are its strongest
products, they are obtained from telescopes that exist, and they are needed
by everything downstream. The pole enters the Yarkovsky drift through
\(\cos\gamma\); the period enters the YORP torque and the lightcurve
amplitude that goes with it.

It is good for identifying the objects worth looking at. A lightcurve is
cheap and it is fast, and an object whose lightcurve shows a large amplitude
and a short period is an object whose shape matters.

It is good for finding the objects that a close approach will change. The
spin state is the observable through which a close encounter reveals the
mass distribution, and the lightcurve is how the spin state is measured.

What it is not good for is the volume. And the volume is what the energy
chain is built from.

\begin{keytakeaways}
\item The standard lightcurve inversion returns the convex hull of the
consistent set, so the convexity assumption is made silently.
\item The field has produced roughly 25,000 photometric models against fewer
than fifteen fly-by shapes and a projected hundred interferometric ones.
\item The Palma silhouette shows real bodies departing from the convex ideal
by tens of per cent, and the hull error is always one-sided.
\item The lightcurve's strongest products are the period and the pole, and
its weakest is the volume, which is the quantity the energy chain needs.
\end{keytakeaways}
